% =====================================================================
%  Anatomía de un LLM — Fase 0 — Módulo 0.3
%  Números, álgebra y sumatorios
%  Compilar: latexmk -r ../latexmkrc P03-algebra.tex
% =====================================================================
\documentclass[11pt,a4paper,twoside,openany]{book}
\usepackage{../estilo-apuntes}

\newcommand{\ApuntesCurso}{Anatomía de un LLM}
\newcommand{\ApuntesModulo}{Módulo 0.3}

% Anclas de hiperenlace únicas por tema para los enunciados
\makeatletter
\@for\nombre:=definicion,teorema,proposicion,lema,corolario,ejemplo,observacion\do{%
  \expandafter\xdef\csname theH\nombre\endcsname{\noexpand\thechapter.\noexpand\arabic{definicion}}}
\makeatother

\newcommand{\N}{\mathbb{N}}
\newcommand{\Z}{\mathbb{Z}}
\newcommand{\Q}{\mathbb{Q}}
\newcommand{\fl}{\operatorname{fl}}
\newcommand{\sgn}{\operatorname{sgn}}
\newcommand{\bits}[1]{\texttt{#1}}

\hypersetup{pdftitle={Módulo 0.3. Números, álgebra y sumatorios (resumen)},pdfauthor={Anatomía de un LLM}}

\makeatletter
\newlabel{ej:periodico}{{1.2}{1}{}{ejemplo.1.2}{}}
\newlabel{ej:periodico@cref}{{[ejemplo][2][]1.2}{[1][1][]1}}
\newlabel{eje:1-clasificar}{{1.1}{4}{}{ejercicio.1.1}{}}
\newlabel{eje:1-clasificar@cref}{{[ejercicio][1][1]1.1}{[1][4][]4}}
\newlabel{eje:1-axiomas}{{1.2}{4}{}{ejercicio.1.2}{}}
\newlabel{eje:1-axiomas@cref}{{[ejercicio][2][1]1.2}{[1][4][]4}}
\newlabel{eje:1-vf}{{1.3}{4}{}{ejercicio.1.3}{}}
\newlabel{eje:1-vf@cref}{{[ejercicio][3][1]1.3}{[1][4][]4}}
\newlabel{eje:1-fracciones}{{1.4}{4}{}{ejercicio.1.4}{}}
\newlabel{eje:1-fracciones@cref}{{[ejercicio][4][1]1.4}{[1][4][]4}}
\newlabel{eje:1-potencias}{{1.5}{4}{}{ejercicio.1.5}{}}
\newlabel{eje:1-potencias@cref}{{[ejercicio][5][1]1.5}{[1][4][]4}}
\newlabel{eje:1-racionalizar}{{1.6}{4}{}{ejercicio.1.6}{}}
\newlabel{eje:1-racionalizar@cref}{{[ejercicio][6][1]1.6}{[1][4][]4}}
\newlabel{eje:1-ordenar}{{1.7}{4}{}{ejercicio.1.7}{}}
\newlabel{eje:1-ordenar@cref}{{[ejercicio][7][1]1.7}{[1][4][]4}}
\newlabel{eje:1-cientifica}{{1.8}{4}{}{ejercicio.1.8}{}}
\newlabel{eje:1-cientifica@cref}{{[ejercicio][8][1]1.8}{[1][4][]4}}
\newlabel{eje:1-intervalos}{{1.9}{4}{}{ejercicio.1.9}{}}
\newlabel{eje:1-intervalos@cref}{{[ejercicio][9][1]1.9}{[1][4][]4}}
\newlabel{eje:1-ecuacion-producto}{{1.10}{5}{}{ejercicio.1.10}{}}
\newlabel{eje:1-ecuacion-producto@cref}{{[ejercicio][10][1]1.10}{[1][4][]5}}
\newlabel{eje:1-quebrados}{{1.11}{5}{}{ejercicio.1.11}{}}
\newlabel{eje:1-quebrados@cref}{{[ejercicio][11][1]1.11}{[1][5][]5}}
\newlabel{eje:1-negativo}{{1.12}{5}{}{ejercicio.1.12}{}}
\newlabel{eje:1-negativo@cref}{{[ejercicio][12][1]1.12}{[1][5][]5}}
\newlabel{eje:1-potencia-n}{{1.13}{5}{}{ejercicio.1.13}{}}
\newlabel{eje:1-potencia-n@cref}{{[ejercicio][13][1]1.13}{[1][5][]5}}
\newlabel{eje:1-raiz3}{{1.14}{5}{}{ejercicio.1.14}{}}
\newlabel{eje:1-raiz3@cref}{{[ejercicio][14][1]1.14}{[1][5][]5}}
\newlabel{eje:1-epsilon}{{1.15}{5}{}{ejercicio.1.15}{}}
\newlabel{eje:1-epsilon@cref}{{[ejercicio][15][1]1.15}{[1][5][]5}}
\newlabel{ej:abs-casos}{{2.5}{7}{}{ejemplo.2.5}{}}
\newlabel{ej:abs-casos@cref}{{[ejemplo][5][]2.5}{[1][7][]7}}
\newlabel{ej:amgm-optim}{{2.9}{8}{}{ejemplo.2.9}{}}
\newlabel{ej:amgm-optim@cref}{{[ejemplo][9][]2.9}{[1][7][]8}}
\newlabel{eje:2-calcular}{{2.1}{9}{}{ejercicio.2.1}{}}
\newlabel{eje:2-calcular@cref}{{[ejercicio][1][2]2.1}{[1][8][]9}}
\newlabel{eje:2-vf}{{2.2}{9}{}{ejercicio.2.2}{}}
\newlabel{eje:2-vf@cref}{{[ejercicio][2][2]2.2}{[1][9][]9}}
\newlabel{eje:2-traducir}{{2.3}{9}{}{ejercicio.2.3}{}}
\newlabel{eje:2-traducir@cref}{{[ejercicio][3][2]2.3}{[1][9][]9}}
\newlabel{eje:2-inec-simple}{{2.4}{9}{}{ejercicio.2.4}{}}
\newlabel{eje:2-inec-simple@cref}{{[ejercicio][4][2]2.4}{[1][9][]9}}
\newlabel{eje:2-inec-casos}{{2.5}{9}{}{ejercicio.2.5}{}}
\newlabel{eje:2-inec-casos@cref}{{[ejercicio][5][2]2.5}{[1][9][]9}}
\newlabel{eje:2-inec-suma}{{2.6}{9}{}{ejercicio.2.6}{}}
\newlabel{eje:2-inec-suma@cref}{{[ejercicio][6][2]2.6}{[1][9][]9}}
\newlabel{eje:2-amgm-aplic}{{2.7}{9}{}{ejercicio.2.7}{}}
\newlabel{eje:2-amgm-aplic@cref}{{[ejercicio][7][2]2.7}{[1][9][]9}}
\newlabel{eje:2-bernoulli-num}{{2.8}{9}{}{ejercicio.2.8}{}}
\newlabel{eje:2-bernoulli-num@cref}{{[ejercicio][8][2]2.8}{[1][9][]9}}
\newlabel{eje:2-cs-num}{{2.9}{9}{}{ejercicio.2.9}{}}
\newlabel{eje:2-cs-num@cref}{{[ejercicio][9][2]2.9}{[1][9][]9}}
\newlabel{eje:2-triangular-n}{{2.10}{9}{}{ejercicio.2.10}{}}
\newlabel{eje:2-triangular-n@cref}{{[ejercicio][10][2]2.10}{[1][9][]9}}
\newlabel{eje:2-maxmin}{{2.11}{9}{}{ejercicio.2.11}{}}
\newlabel{eje:2-maxmin@cref}{{[ejercicio][11][2]2.11}{[1][9][]9}}
\newlabel{eje:2-cadena-medias}{{2.12}{9}{}{ejercicio.2.12}{}}
\newlabel{eje:2-cadena-medias@cref}{{[ejercicio][12][2]2.12}{[1][9][]9}}
\newlabel{eje:2-media-armonica}{{2.13}{9}{}{ejercicio.2.13}{}}
\newlabel{eje:2-media-armonica@cref}{{[ejercicio][13][2]2.13}{[1][9][]9}}
\newlabel{eje:2-bernoulli-fuerte}{{2.14}{10}{}{ejercicio.2.14}{}}
\newlabel{eje:2-bernoulli-fuerte@cref}{{[ejercicio][14][2]2.14}{[1][9][]10}}
\newlabel{eje:2-suma-inversos}{{2.15}{10}{}{ejercicio.2.15}{}}
\newlabel{eje:2-suma-inversos@cref}{{[ejercicio][15][2]2.15}{[1][10][]10}}
\newlabel{ej:ruffini}{{3.4}{12}{}{ejemplo.3.4}{}}
\newlabel{ej:ruffini@cref}{{[ejemplo][4][]3.4}{[1][11][]12}}
\newlabel{ej:tabla-signos}{{3.10}{13}{}{ejemplo.3.10}{}}
\newlabel{ej:tabla-signos@cref}{{[ejemplo][10][]3.10}{[1][13][]13}}
\newlabel{ej:gauss3}{{3.12}{13}{}{ejemplo.3.12}{}}
\newlabel{ej:gauss3@cref}{{[ejemplo][12][]3.12}{[1][13][]13}}
\newlabel{eje:3-basico}{{3.1}{14}{}{ejercicio.3.1}{}}
\newlabel{eje:3-basico@cref}{{[ejercicio][1][3]3.1}{[1][14][]14}}
\newlabel{eje:3-discriminante}{{3.2}{14}{}{ejercicio.3.2}{}}
\newlabel{eje:3-discriminante@cref}{{[ejercicio][2][3]3.2}{[1][14][]14}}
\newlabel{eje:3-division}{{3.3}{14}{}{ejercicio.3.3}{}}
\newlabel{eje:3-division@cref}{{[ejercicio][3][3]3.3}{[1][14][]14}}
\newlabel{eje:3-ruffini}{{3.4}{14}{}{ejercicio.3.4}{}}
\newlabel{eje:3-ruffini@cref}{{[ejercicio][4][3]3.4}{[1][14][]14}}
\newlabel{eje:3-resto}{{3.5}{14}{}{ejercicio.3.5}{}}
\newlabel{eje:3-resto@cref}{{[ejercicio][5][3]3.5}{[1][14][]14}}
\newlabel{eje:3-racionales}{{3.6}{14}{}{ejercicio.3.6}{}}
\newlabel{eje:3-racionales@cref}{{[ejercicio][6][3]3.6}{[1][14][]14}}
\newlabel{eje:3-vieta}{{3.7}{15}{}{ejercicio.3.7}{}}
\newlabel{eje:3-vieta@cref}{{[ejercicio][7][3]3.7}{[1][14][]15}}
\newlabel{eje:3-ecuaciones}{{3.8}{15}{}{ejercicio.3.8}{}}
\newlabel{eje:3-ecuaciones@cref}{{[ejercicio][8][3]3.8}{[1][15][]15}}
\newlabel{eje:3-inecuaciones}{{3.9}{15}{}{ejercicio.3.9}{}}
\newlabel{eje:3-inecuaciones@cref}{{[ejercicio][9][3]3.9}{[1][15][]15}}
\newlabel{eje:3-sistemas}{{3.10}{15}{}{ejercicio.3.10}{}}
\newlabel{eje:3-sistemas@cref}{{[ejercicio][10][3]3.10}{[1][15][]15}}
\newlabel{eje:3-identidad}{{3.11}{15}{}{ejercicio.3.11}{}}
\newlabel{eje:3-identidad@cref}{{[ejercicio][11][3]3.11}{[1][15][]15}}
\newlabel{eje:3-diferencia-potencias}{{3.12}{15}{}{ejercicio.3.12}{}}
\newlabel{eje:3-diferencia-potencias@cref}{{[ejercicio][12][3]3.12}{[1][15][]15}}
\newlabel{eje:3-impares}{{3.13}{15}{}{ejercicio.3.13}{}}
\newlabel{eje:3-impares@cref}{{[ejercicio][13][3]3.13}{[1][15][]15}}
\newlabel{eje:3-det-cero}{{3.14}{15}{}{ejercicio.3.14}{}}
\newlabel{eje:3-det-cero@cref}{{[ejercicio][14][3]3.14}{[1][15][]15}}
\newlabel{eje:3-cuadratica-positiva}{{3.15}{15}{}{ejercicio.3.15}{}}
\newlabel{eje:3-cuadratica-positiva@cref}{{[ejercicio][15][3]3.15}{[1][15][]15}}
\newlabel{ej:telescopica}{{4.7}{17}{}{ejemplo.4.7}{}}
\newlabel{ej:telescopica@cref}{{[ejemplo][7][]4.7}{[1][17][]17}}
\newlabel{ej:ema}{{4.9}{18}{Medias exponenciales}{ejemplo.4.9}{}}
\newlabel{ej:ema@cref}{{[ejemplo][9][]4.9}{[1][17][]18}}
\newlabel{ej:causal}{{4.11}{18}{Coste de la atención causal}{ejemplo.4.11}{}}
\newlabel{ej:causal@cref}{{[ejemplo][11][]4.11}{[1][18][]18}}
\newlabel{eje:4-notacion}{{4.1}{19}{}{ejercicio.4.1}{}}
\newlabel{eje:4-notacion@cref}{{[ejercicio][1][4]4.1}{[1][19][]19}}
\newlabel{eje:4-calcular}{{4.2}{19}{}{ejercicio.4.2}{}}
\newlabel{eje:4-calcular@cref}{{[ejercicio][2][4]4.2}{[1][19][]19}}
\newlabel{eje:4-errores}{{4.3}{20}{}{ejercicio.4.3}{}}
\newlabel{eje:4-errores@cref}{{[ejercicio][3][4]4.3}{[1][20][]20}}
\newlabel{eje:4-aritmeticas}{{4.4}{20}{}{ejercicio.4.4}{}}
\newlabel{eje:4-aritmeticas@cref}{{[ejercicio][4][4]4.4}{[1][20][]20}}
\newlabel{eje:4-geometricas}{{4.5}{20}{}{ejercicio.4.5}{}}
\newlabel{eje:4-geometricas@cref}{{[ejercicio][5][4]4.5}{[1][20][]20}}
\newlabel{eje:4-telescopicas}{{4.6}{20}{}{ejercicio.4.6}{}}
\newlabel{eje:4-telescopicas@cref}{{[ejercicio][6][4]4.6}{[1][20][]20}}
\newlabel{eje:4-indice}{{4.7}{20}{}{ejercicio.4.7}{}}
\newlabel{eje:4-indice@cref}{{[ejercicio][7][4]4.7}{[1][20][]20}}
\newlabel{eje:4-dobles}{{4.8}{20}{}{ejercicio.4.8}{}}
\newlabel{eje:4-dobles@cref}{{[ejercicio][8][4]4.8}{[1][20][]20}}
\newlabel{eje:4-binomio}{{4.9}{20}{}{ejercicio.4.9}{}}
\newlabel{eje:4-binomio@cref}{{[ejercicio][9][4]4.9}{[1][20][]20}}
\newlabel{eje:4-productos}{{4.10}{20}{}{ejercicio.4.10}{}}
\newlabel{eje:4-productos@cref}{{[ejercicio][10][4]4.10}{[1][20][]20}}
\newlabel{eje:4-cubos}{{4.11}{20}{}{ejercicio.4.11}{}}
\newlabel{eje:4-cubos@cref}{{[ejercicio][11][4]4.11}{[1][20][]20}}
\newlabel{eje:4-perturbacion}{{4.12}{20}{}{ejercicio.4.12}{}}
\newlabel{eje:4-perturbacion@cref}{{[ejercicio][12][4]4.12}{[1][20][]20}}
\newlabel{eje:4-kbinom}{{4.13}{20}{}{ejercicio.4.13}{}}
\newlabel{eje:4-kbinom@cref}{{[ejercicio][13][4]4.13}{[1][20][]20}}
\newlabel{eje:4-basilea}{{4.14}{21}{}{ejercicio.4.14}{}}
\newlabel{eje:4-basilea@cref}{{[ejercicio][14][4]4.14}{[1][20][]21}}
\newlabel{eje:4-varianza}{{4.15}{21}{}{ejercicio.4.15}{}}
\newlabel{eje:4-varianza@cref}{{[ejercicio][15][4]4.15}{[1][21][]21}}
\newlabel{ej:ecuacion-exp}{{5.6}{23}{}{ejemplo.5.6}{}}
\newlabel{ej:ecuacion-exp@cref}{{[ejemplo][6][]5.6}{[1][23][]23}}
\newlabel{ej:ecuacion-log}{{5.7}{23}{}{ejemplo.5.7}{}}
\newlabel{ej:ecuacion-log@cref}{{[ejemplo][7][]5.7}{[1][23][]23}}
\newlabel{ej:lse-estable}{{5.10}{24}{}{ejemplo.5.10}{}}
\newlabel{ej:lse-estable@cref}{{[ejemplo][10][]5.10}{[1][24][]24}}
\newlabel{ej:underflow}{{5.12}{24}{Productos de probabilidades}{ejemplo.5.12}{}}
\newlabel{ej:underflow@cref}{{[ejemplo][12][]5.12}{[1][24][]24}}
\newlabel{eje:5-calcular}{{5.1}{25}{}{ejercicio.5.1}{}}
\newlabel{eje:5-calcular@cref}{{[ejercicio][1][5]5.1}{[1][25][]25}}
\newlabel{eje:5-vf}{{5.2}{25}{}{ejercicio.5.2}{}}
\newlabel{eje:5-vf@cref}{{[ejercicio][2][5]5.2}{[1][25][]25}}
\newlabel{eje:5-unidades}{{5.3}{25}{}{ejercicio.5.3}{}}
\newlabel{eje:5-unidades@cref}{{[ejercicio][3][5]5.3}{[1][25][]25}}
\newlabel{eje:5-simplificar}{{5.4}{25}{}{ejercicio.5.4}{}}
\newlabel{eje:5-simplificar@cref}{{[ejercicio][4][5]5.4}{[1][25][]25}}
\newlabel{eje:5-exponenciales}{{5.5}{25}{}{ejercicio.5.5}{}}
\newlabel{eje:5-exponenciales@cref}{{[ejercicio][5][5]5.5}{[1][25][]25}}
\newlabel{eje:5-logaritmicas}{{5.6}{25}{}{ejercicio.5.6}{}}
\newlabel{eje:5-logaritmicas@cref}{{[ejercicio][6][5]5.6}{[1][25][]25}}
\newlabel{eje:5-inecuaciones}{{5.7}{25}{}{ejercicio.5.7}{}}
\newlabel{eje:5-inecuaciones@cref}{{[ejercicio][7][5]5.7}{[1][25][]25}}
\newlabel{eje:5-lse}{{5.8}{25}{}{ejercicio.5.8}{}}
\newlabel{eje:5-lse@cref}{{[ejercicio][8][5]5.8}{[1][25][]25}}
\newlabel{eje:5-perplejidad}{{5.9}{25}{}{ejercicio.5.9}{}}
\newlabel{eje:5-perplejidad@cref}{{[ejercicio][9][5]5.9}{[1][25][]25}}
\newlabel{eje:5-ley-potencia}{{5.10}{25}{}{ejercicio.5.10}{}}
\newlabel{eje:5-ley-potencia@cref}{{[ejercicio][10][5]5.10}{[1][25][]25}}
\newlabel{eje:5-log23}{{5.11}{25}{}{ejercicio.5.11}{}}
\newlabel{eje:5-log23@cref}{{[ejercicio][11][5]5.11}{[1][25][]25}}
\newlabel{eje:5-temperatura}{{5.12}{25}{}{ejercicio.5.12}{}}
\newlabel{eje:5-temperatura@cref}{{[ejercicio][12][5]5.12}{[1][25][]25}}
\newlabel{eje:5-lse-media}{{5.13}{26}{}{ejercicio.5.13}{}}
\newlabel{eje:5-lse-media@cref}{{[ejercicio][13][5]5.13}{[1][25][]26}}
\newlabel{eje:5-gibbs}{{5.14}{26}{}{ejercicio.5.14}{}}
\newlabel{eje:5-gibbs@cref}{{[ejercicio][14][5]5.14}{[1][26][]26}}
\newlabel{eje:5-e-creciente}{{5.15}{26}{}{ejercicio.5.15}{}}
\newlabel{eje:5-e-creciente@cref}{{[ejercicio][15][5]5.15}{[1][26][]26}}
\newlabel{ej:conversion}{{6.2}{27}{}{ejemplo.6.2}{}}
\newlabel{ej:conversion@cref}{{[ejemplo][2][]6.2}{[1][27][]27}}
\newlabel{ej:c2}{{6.7}{29}{}{ejemplo.6.7}{}}
\newlabel{ej:c2@cref}{{[ejemplo][7][]6.7}{[1][29][]29}}
\newlabel{eje:6-a-decimal}{{6.1}{30}{}{ejercicio.6.1}{}}
\newlabel{eje:6-a-decimal@cref}{{[ejercicio][1][6]6.1}{[1][30][]30}}
\newlabel{eje:6-contar}{{6.2}{30}{}{ejercicio.6.2}{}}
\newlabel{eje:6-contar@cref}{{[ejercicio][2][6]6.2}{[1][30][]30}}
\newlabel{eje:6-vf}{{6.3}{30}{}{ejercicio.6.3}{}}
\newlabel{eje:6-vf@cref}{{[ejercicio][3][6]6.3}{[1][30][]30}}
\newlabel{eje:6-a-binario}{{6.4}{30}{}{ejercicio.6.4}{}}
\newlabel{eje:6-a-binario@cref}{{[ejercicio][4][6]6.4}{[1][30][]30}}
\newlabel{eje:6-fracciones}{{6.5}{30}{}{ejercicio.6.5}{}}
\newlabel{eje:6-fracciones@cref}{{[ejercicio][5][6]6.5}{[1][30][]30}}
\newlabel{eje:6-c2}{{6.6}{30}{}{ejercicio.6.6}{}}
\newlabel{eje:6-c2@cref}{{[ejercicio][6][6]6.6}{[1][30][]30}}
\newlabel{eje:6-utf8}{{6.7}{30}{}{ejercicio.6.7}{}}
\newlabel{eje:6-utf8@cref}{{[ejercicio][7][6]6.7}{[1][30][]30}}
\newlabel{eje:6-digitos}{{6.8}{31}{}{ejercicio.6.8}{}}
\newlabel{eje:6-digitos@cref}{{[ejercicio][8][6]6.8}{[1][30][]31}}
\newlabel{eje:6-decodificar}{{6.9}{31}{}{ejercicio.6.9}{}}
\newlabel{eje:6-decodificar@cref}{{[ejercicio][9][6]6.9}{[1][31][]31}}
\newlabel{eje:6-espaciado}{{6.10}{31}{}{ejercicio.6.10}{}}
\newlabel{eje:6-espaciado@cref}{{[ejercicio][10][6]6.10}{[1][31][]31}}
\newlabel{eje:6-memoria}{{6.11}{31}{}{ejercicio.6.11}{}}
\newlabel{eje:6-memoria@cref}{{[ejercicio][11][6]6.11}{[1][31][]31}}
\newlabel{eje:6-divisibilidad}{{6.12}{31}{}{ejercicio.6.12}{}}
\newlabel{eje:6-divisibilidad@cref}{{[ejercicio][12][6]6.12}{[1][31][]31}}
\newlabel{eje:6-extension}{{6.13}{31}{}{ejercicio.6.13}{}}
\newlabel{eje:6-extension@cref}{{[ejercicio][13][6]6.13}{[1][31][]31}}
\newlabel{eje:6-decimal-finito}{{6.14}{31}{}{ejercicio.6.14}{}}
\newlabel{eje:6-decimal-finito@cref}{{[ejercicio][14][6]6.14}{[1][31][]31}}
\newlabel{eje:6-desbordamiento}{{6.15}{31}{}{ejercicio.6.15}{}}
\newlabel{eje:6-desbordamiento@cref}{{[ejercicio][15][6]6.15}{[1][31][]31}}
\newlabel{eje:6-redondeo}{{6.16}{31}{}{ejercicio.6.16}{}}
\newlabel{eje:6-redondeo@cref}{{[ejercicio][16][6]6.16}{[1][31][]31}}
\newlabel{ap:soluciones}{{A}{32}{Soluciones de los ejercicios}{appendix.A}{}}
\newlabel{ap:soluciones@cref}{{[appendix][1][2147483647]A}{[1][32][]32}}
\newlabel{sol:eje:1-clasificar}{{A.1}{32}{Tema 1}{section*.3}{}}
\newlabel{sol:eje:1-clasificar@cref}{{[subappendix][1][2147483647,1]A.1}{[1][32][]32}}
\newlabel{sol:eje:1-axiomas}{{A.1}{32}{Tema 1}{section*.4}{}}
\newlabel{sol:eje:1-axiomas@cref}{{[subappendix][1][2147483647,1]A.1}{[1][32][]32}}
\newlabel{sol:eje:1-vf}{{A.1}{33}{Tema 1}{section*.5}{}}
\newlabel{sol:eje:1-vf@cref}{{[subappendix][1][2147483647,1]A.1}{[1][33][]33}}
\newlabel{sol:eje:1-fracciones}{{A.1}{33}{Tema 1}{section*.6}{}}
\newlabel{sol:eje:1-fracciones@cref}{{[subappendix][1][2147483647,1]A.1}{[1][33][]33}}
\newlabel{sol:eje:1-potencias}{{A.1}{33}{Tema 1}{section*.7}{}}
\newlabel{sol:eje:1-potencias@cref}{{[subappendix][1][2147483647,1]A.1}{[1][33][]33}}
\newlabel{sol:eje:1-racionalizar}{{A.1}{33}{Tema 1}{section*.8}{}}
\newlabel{sol:eje:1-racionalizar@cref}{{[subappendix][1][2147483647,1]A.1}{[1][33][]33}}
\newlabel{sol:eje:1-ordenar}{{A.1}{34}{Tema 1}{section*.9}{}}
\newlabel{sol:eje:1-ordenar@cref}{{[subappendix][1][2147483647,1]A.1}{[1][34][]34}}
\newlabel{sol:eje:1-cientifica}{{A.1}{34}{Tema 1}{section*.10}{}}
\newlabel{sol:eje:1-cientifica@cref}{{[subappendix][1][2147483647,1]A.1}{[1][34][]34}}
\newlabel{sol:eje:1-intervalos}{{A.1}{34}{Tema 1}{section*.11}{}}
\newlabel{sol:eje:1-intervalos@cref}{{[subappendix][1][2147483647,1]A.1}{[1][34][]34}}
\newlabel{sol:eje:1-ecuacion-producto}{{A.1}{35}{Tema 1}{section*.12}{}}
\newlabel{sol:eje:1-ecuacion-producto@cref}{{[subappendix][1][2147483647,1]A.1}{[1][35][]35}}
\newlabel{sol:eje:1-quebrados}{{A.1}{35}{Tema 1}{section*.13}{}}
\newlabel{sol:eje:1-quebrados@cref}{{[subappendix][1][2147483647,1]A.1}{[1][35][]35}}
\newlabel{sol:eje:1-negativo}{{A.1}{35}{Tema 1}{section*.14}{}}
\newlabel{sol:eje:1-negativo@cref}{{[subappendix][1][2147483647,1]A.1}{[1][35][]35}}
\newlabel{sol:eje:1-potencia-n}{{A.1}{36}{Tema 1}{section*.15}{}}
\newlabel{sol:eje:1-potencia-n@cref}{{[subappendix][1][2147483647,1]A.1}{[1][36][]36}}
\newlabel{sol:eje:1-raiz3}{{A.1}{36}{Tema 1}{section*.16}{}}
\newlabel{sol:eje:1-raiz3@cref}{{[subappendix][1][2147483647,1]A.1}{[1][36][]36}}
\newlabel{sol:eje:1-epsilon}{{A.1}{36}{Tema 1}{section*.17}{}}
\newlabel{sol:eje:1-epsilon@cref}{{[subappendix][1][2147483647,1]A.1}{[1][36][]36}}
\newlabel{sol:eje:2-calcular}{{A.2}{37}{Tema 2}{section*.18}{}}
\newlabel{sol:eje:2-calcular@cref}{{[subappendix][2][2147483647,1]A.2}{[1][37][]37}}
\newlabel{sol:eje:2-vf}{{A.2}{37}{Tema 2}{section*.19}{}}
\newlabel{sol:eje:2-vf@cref}{{[subappendix][2][2147483647,1]A.2}{[1][37][]37}}
\newlabel{sol:eje:2-traducir}{{A.2}{37}{Tema 2}{section*.20}{}}
\newlabel{sol:eje:2-traducir@cref}{{[subappendix][2][2147483647,1]A.2}{[1][37][]37}}
\newlabel{sol:eje:2-inec-simple}{{A.2}{38}{Tema 2}{section*.21}{}}
\newlabel{sol:eje:2-inec-simple@cref}{{[subappendix][2][2147483647,1]A.2}{[1][38][]38}}
\newlabel{sol:eje:2-inec-casos}{{A.2}{38}{Tema 2}{section*.22}{}}
\newlabel{sol:eje:2-inec-casos@cref}{{[subappendix][2][2147483647,1]A.2}{[1][38][]38}}
\newlabel{sol:eje:2-inec-suma}{{A.2}{38}{Tema 2}{section*.23}{}}
\newlabel{sol:eje:2-inec-suma@cref}{{[subappendix][2][2147483647,1]A.2}{[1][38][]38}}
\newlabel{sol:eje:2-amgm-aplic}{{A.2}{39}{Tema 2}{section*.24}{}}
\newlabel{sol:eje:2-amgm-aplic@cref}{{[subappendix][2][2147483647,1]A.2}{[1][38][]39}}
\newlabel{sol:eje:2-bernoulli-num}{{A.2}{39}{Tema 2}{section*.25}{}}
\newlabel{sol:eje:2-bernoulli-num@cref}{{[subappendix][2][2147483647,1]A.2}{[1][39][]39}}
\newlabel{sol:eje:2-cs-num}{{A.2}{39}{Tema 2}{section*.26}{}}
\newlabel{sol:eje:2-cs-num@cref}{{[subappendix][2][2147483647,1]A.2}{[1][39][]39}}
\newlabel{sol:eje:2-triangular-n}{{A.2}{39}{Tema 2}{section*.27}{}}
\newlabel{sol:eje:2-triangular-n@cref}{{[subappendix][2][2147483647,1]A.2}{[1][39][]39}}
\newlabel{sol:eje:2-maxmin}{{A.2}{40}{Tema 2}{section*.28}{}}
\newlabel{sol:eje:2-maxmin@cref}{{[subappendix][2][2147483647,1]A.2}{[1][40][]40}}
\newlabel{sol:eje:2-cadena-medias}{{A.2}{40}{Tema 2}{section*.29}{}}
\newlabel{sol:eje:2-cadena-medias@cref}{{[subappendix][2][2147483647,1]A.2}{[1][40][]40}}
\newlabel{sol:eje:2-media-armonica}{{A.2}{40}{Tema 2}{section*.30}{}}
\newlabel{sol:eje:2-media-armonica@cref}{{[subappendix][2][2147483647,1]A.2}{[1][40][]40}}
\newlabel{sol:eje:2-bernoulli-fuerte}{{A.2}{41}{Tema 2}{section*.31}{}}
\newlabel{sol:eje:2-bernoulli-fuerte@cref}{{[subappendix][2][2147483647,1]A.2}{[1][40][]41}}
\newlabel{sol:eje:2-suma-inversos}{{A.2}{41}{Tema 2}{section*.32}{}}
\newlabel{sol:eje:2-suma-inversos@cref}{{[subappendix][2][2147483647,1]A.2}{[1][41][]41}}
\newlabel{sol:eje:3-basico}{{A.3}{41}{Tema 3}{section*.33}{}}
\newlabel{sol:eje:3-basico@cref}{{[subappendix][3][2147483647,1]A.3}{[1][41][]41}}
\newlabel{sol:eje:3-discriminante}{{A.3}{42}{Tema 3}{section*.34}{}}
\newlabel{sol:eje:3-discriminante@cref}{{[subappendix][3][2147483647,1]A.3}{[1][42][]42}}
\newlabel{sol:eje:3-division}{{A.3}{42}{Tema 3}{section*.35}{}}
\newlabel{sol:eje:3-division@cref}{{[subappendix][3][2147483647,1]A.3}{[1][42][]42}}
\newlabel{sol:eje:3-ruffini}{{A.3}{42}{Tema 3}{section*.36}{}}
\newlabel{sol:eje:3-ruffini@cref}{{[subappendix][3][2147483647,1]A.3}{[1][42][]42}}
\newlabel{sol:eje:3-resto}{{A.3}{42}{Tema 3}{section*.37}{}}
\newlabel{sol:eje:3-resto@cref}{{[subappendix][3][2147483647,1]A.3}{[1][42][]42}}
\newlabel{sol:eje:3-racionales}{{A.3}{43}{Tema 3}{section*.38}{}}
\newlabel{sol:eje:3-racionales@cref}{{[subappendix][3][2147483647,1]A.3}{[1][43][]43}}
\newlabel{sol:eje:3-vieta}{{A.3}{43}{Tema 3}{section*.39}{}}
\newlabel{sol:eje:3-vieta@cref}{{[subappendix][3][2147483647,1]A.3}{[1][43][]43}}
\newlabel{sol:eje:3-ecuaciones}{{A.3}{43}{Tema 3}{section*.40}{}}
\newlabel{sol:eje:3-ecuaciones@cref}{{[subappendix][3][2147483647,1]A.3}{[1][43][]43}}
\newlabel{sol:eje:3-inecuaciones}{{A.3}{44}{Tema 3}{section*.41}{}}
\newlabel{sol:eje:3-inecuaciones@cref}{{[subappendix][3][2147483647,1]A.3}{[1][43][]44}}
\newlabel{sol:eje:3-sistemas}{{A.3}{44}{Tema 3}{section*.42}{}}
\newlabel{sol:eje:3-sistemas@cref}{{[subappendix][3][2147483647,1]A.3}{[1][44][]44}}
\newlabel{sol:eje:3-identidad}{{A.3}{44}{Tema 3}{section*.43}{}}
\newlabel{sol:eje:3-identidad@cref}{{[subappendix][3][2147483647,1]A.3}{[1][44][]44}}
\newlabel{sol:eje:3-diferencia-potencias}{{A.3}{44}{Tema 3}{section*.44}{}}
\newlabel{sol:eje:3-diferencia-potencias@cref}{{[subappendix][3][2147483647,1]A.3}{[1][44][]44}}
\newlabel{sol:eje:3-impares}{{A.3}{45}{Tema 3}{section*.45}{}}
\newlabel{sol:eje:3-impares@cref}{{[subappendix][3][2147483647,1]A.3}{[1][45][]45}}
\newlabel{sol:eje:3-det-cero}{{A.3}{45}{Tema 3}{section*.46}{}}
\newlabel{sol:eje:3-det-cero@cref}{{[subappendix][3][2147483647,1]A.3}{[1][45][]45}}
\newlabel{sol:eje:3-cuadratica-positiva}{{A.3}{46}{Tema 3}{section*.47}{}}
\newlabel{sol:eje:3-cuadratica-positiva@cref}{{[subappendix][3][2147483647,1]A.3}{[1][45][]46}}
\newlabel{sol:eje:4-notacion}{{A.4}{46}{Tema 4}{section*.48}{}}
\newlabel{sol:eje:4-notacion@cref}{{[subappendix][4][2147483647,1]A.4}{[1][46][]46}}
\newlabel{sol:eje:4-calcular}{{A.4}{46}{Tema 4}{section*.49}{}}
\newlabel{sol:eje:4-calcular@cref}{{[subappendix][4][2147483647,1]A.4}{[1][46][]46}}
\newlabel{sol:eje:4-errores}{{A.4}{46}{Tema 4}{section*.50}{}}
\newlabel{sol:eje:4-errores@cref}{{[subappendix][4][2147483647,1]A.4}{[1][46][]46}}
\newlabel{sol:eje:4-aritmeticas}{{A.4}{47}{Tema 4}{section*.51}{}}
\newlabel{sol:eje:4-aritmeticas@cref}{{[subappendix][4][2147483647,1]A.4}{[1][47][]47}}
\newlabel{sol:eje:4-geometricas}{{A.4}{47}{Tema 4}{section*.52}{}}
\newlabel{sol:eje:4-geometricas@cref}{{[subappendix][4][2147483647,1]A.4}{[1][47][]47}}
\newlabel{sol:eje:4-telescopicas}{{A.4}{47}{Tema 4}{section*.53}{}}
\newlabel{sol:eje:4-telescopicas@cref}{{[subappendix][4][2147483647,1]A.4}{[1][47][]47}}
\newlabel{sol:eje:4-indice}{{A.4}{48}{Tema 4}{section*.54}{}}
\newlabel{sol:eje:4-indice@cref}{{[subappendix][4][2147483647,1]A.4}{[1][47][]48}}
\newlabel{sol:eje:4-dobles}{{A.4}{48}{Tema 4}{section*.55}{}}
\newlabel{sol:eje:4-dobles@cref}{{[subappendix][4][2147483647,1]A.4}{[1][48][]48}}
\newlabel{sol:eje:4-binomio}{{A.4}{48}{Tema 4}{section*.56}{}}
\newlabel{sol:eje:4-binomio@cref}{{[subappendix][4][2147483647,1]A.4}{[1][48][]48}}
\newlabel{sol:eje:4-productos}{{A.4}{48}{Tema 4}{section*.57}{}}
\newlabel{sol:eje:4-productos@cref}{{[subappendix][4][2147483647,1]A.4}{[1][48][]48}}
\newlabel{sol:eje:4-cubos}{{A.4}{49}{Tema 4}{section*.58}{}}
\newlabel{sol:eje:4-cubos@cref}{{[subappendix][4][2147483647,1]A.4}{[1][49][]49}}
\newlabel{sol:eje:4-perturbacion}{{A.4}{49}{Tema 4}{section*.59}{}}
\newlabel{sol:eje:4-perturbacion@cref}{{[subappendix][4][2147483647,1]A.4}{[1][49][]49}}
\newlabel{sol:eje:4-kbinom}{{A.4}{49}{Tema 4}{section*.60}{}}
\newlabel{sol:eje:4-kbinom@cref}{{[subappendix][4][2147483647,1]A.4}{[1][49][]49}}
\newlabel{sol:eje:4-basilea}{{A.4}{50}{Tema 4}{section*.61}{}}
\newlabel{sol:eje:4-basilea@cref}{{[subappendix][4][2147483647,1]A.4}{[1][50][]50}}
\newlabel{sol:eje:4-varianza}{{A.4}{50}{Tema 4}{section*.62}{}}
\newlabel{sol:eje:4-varianza@cref}{{[subappendix][4][2147483647,1]A.4}{[1][50][]50}}
\newlabel{sol:eje:5-calcular}{{A.5}{50}{Tema 5}{section*.63}{}}
\newlabel{sol:eje:5-calcular@cref}{{[subappendix][5][2147483647,1]A.5}{[1][50][]50}}
\newlabel{sol:eje:5-vf}{{A.5}{51}{Tema 5}{section*.64}{}}
\newlabel{sol:eje:5-vf@cref}{{[subappendix][5][2147483647,1]A.5}{[1][51][]51}}
\newlabel{sol:eje:5-unidades}{{A.5}{51}{Tema 5}{section*.65}{}}
\newlabel{sol:eje:5-unidades@cref}{{[subappendix][5][2147483647,1]A.5}{[1][51][]51}}
\newlabel{sol:eje:5-simplificar}{{A.5}{51}{Tema 5}{section*.66}{}}
\newlabel{sol:eje:5-simplificar@cref}{{[subappendix][5][2147483647,1]A.5}{[1][51][]51}}
\newlabel{sol:eje:5-exponenciales}{{A.5}{52}{Tema 5}{section*.67}{}}
\newlabel{sol:eje:5-exponenciales@cref}{{[subappendix][5][2147483647,1]A.5}{[1][51][]52}}
\newlabel{sol:eje:5-logaritmicas}{{A.5}{52}{Tema 5}{section*.68}{}}
\newlabel{sol:eje:5-logaritmicas@cref}{{[subappendix][5][2147483647,1]A.5}{[1][52][]52}}
\newlabel{sol:eje:5-inecuaciones}{{A.5}{52}{Tema 5}{section*.69}{}}
\newlabel{sol:eje:5-inecuaciones@cref}{{[subappendix][5][2147483647,1]A.5}{[1][52][]52}}
\newlabel{sol:eje:5-lse}{{A.5}{52}{Tema 5}{section*.70}{}}
\newlabel{sol:eje:5-lse@cref}{{[subappendix][5][2147483647,1]A.5}{[1][52][]52}}
\newlabel{sol:eje:5-perplejidad}{{A.5}{53}{Tema 5}{section*.71}{}}
\newlabel{sol:eje:5-perplejidad@cref}{{[subappendix][5][2147483647,1]A.5}{[1][53][]53}}
\newlabel{sol:eje:5-ley-potencia}{{A.5}{53}{Tema 5}{section*.72}{}}
\newlabel{sol:eje:5-ley-potencia@cref}{{[subappendix][5][2147483647,1]A.5}{[1][53][]53}}
\newlabel{sol:eje:5-log23}{{A.5}{53}{Tema 5}{section*.73}{}}
\newlabel{sol:eje:5-log23@cref}{{[subappendix][5][2147483647,1]A.5}{[1][53][]53}}
\newlabel{sol:eje:5-temperatura}{{A.5}{54}{Tema 5}{section*.74}{}}
\newlabel{sol:eje:5-temperatura@cref}{{[subappendix][5][2147483647,1]A.5}{[1][53][]54}}
\newlabel{sol:eje:5-lse-media}{{A.5}{54}{Tema 5}{section*.75}{}}
\newlabel{sol:eje:5-lse-media@cref}{{[subappendix][5][2147483647,1]A.5}{[1][54][]54}}
\newlabel{sol:eje:5-gibbs}{{A.5}{54}{Tema 5}{section*.76}{}}
\newlabel{sol:eje:5-gibbs@cref}{{[subappendix][5][2147483647,1]A.5}{[1][54][]54}}
\newlabel{sol:eje:5-e-creciente}{{A.5}{54}{Tema 5}{section*.77}{}}
\newlabel{sol:eje:5-e-creciente@cref}{{[subappendix][5][2147483647,1]A.5}{[1][54][]54}}
\newlabel{sol:eje:6-a-decimal}{{A.6}{55}{Tema 6}{section*.78}{}}
\newlabel{sol:eje:6-a-decimal@cref}{{[subappendix][6][2147483647,1]A.6}{[1][55][]55}}
\newlabel{sol:eje:6-contar}{{A.6}{55}{Tema 6}{section*.79}{}}
\newlabel{sol:eje:6-contar@cref}{{[subappendix][6][2147483647,1]A.6}{[1][55][]55}}
\newlabel{sol:eje:6-vf}{{A.6}{55}{Tema 6}{section*.80}{}}
\newlabel{sol:eje:6-vf@cref}{{[subappendix][6][2147483647,1]A.6}{[1][55][]55}}
\newlabel{sol:eje:6-a-binario}{{A.6}{56}{Tema 6}{section*.81}{}}
\newlabel{sol:eje:6-a-binario@cref}{{[subappendix][6][2147483647,1]A.6}{[1][56][]56}}
\newlabel{sol:eje:6-fracciones}{{A.6}{56}{Tema 6}{section*.82}{}}
\newlabel{sol:eje:6-fracciones@cref}{{[subappendix][6][2147483647,1]A.6}{[1][56][]56}}
\newlabel{sol:eje:6-c2}{{A.6}{56}{Tema 6}{section*.83}{}}
\newlabel{sol:eje:6-c2@cref}{{[subappendix][6][2147483647,1]A.6}{[1][56][]56}}
\newlabel{sol:eje:6-utf8}{{A.6}{57}{Tema 6}{section*.84}{}}
\newlabel{sol:eje:6-utf8@cref}{{[subappendix][6][2147483647,1]A.6}{[1][56][]57}}
\newlabel{sol:eje:6-digitos}{{A.6}{57}{Tema 6}{section*.85}{}}
\newlabel{sol:eje:6-digitos@cref}{{[subappendix][6][2147483647,1]A.6}{[1][57][]57}}
\newlabel{sol:eje:6-decodificar}{{A.6}{57}{Tema 6}{section*.86}{}}
\newlabel{sol:eje:6-decodificar@cref}{{[subappendix][6][2147483647,1]A.6}{[1][57][]57}}
\newlabel{sol:eje:6-espaciado}{{A.6}{58}{Tema 6}{section*.87}{}}
\newlabel{sol:eje:6-espaciado@cref}{{[subappendix][6][2147483647,1]A.6}{[1][57][]58}}
\newlabel{sol:eje:6-memoria}{{A.6}{58}{Tema 6}{section*.88}{}}
\newlabel{sol:eje:6-memoria@cref}{{[subappendix][6][2147483647,1]A.6}{[1][58][]58}}
\newlabel{sol:eje:6-divisibilidad}{{A.6}{58}{Tema 6}{section*.89}{}}
\newlabel{sol:eje:6-divisibilidad@cref}{{[subappendix][6][2147483647,1]A.6}{[1][58][]58}}
\newlabel{sol:eje:6-extension}{{A.6}{59}{Tema 6}{section*.90}{}}
\newlabel{sol:eje:6-extension@cref}{{[subappendix][6][2147483647,1]A.6}{[1][59][]59}}
\newlabel{sol:eje:6-decimal-finito}{{A.6}{59}{Tema 6}{section*.91}{}}
\newlabel{sol:eje:6-decimal-finito@cref}{{[subappendix][6][2147483647,1]A.6}{[1][59][]59}}
\newlabel{sol:eje:6-desbordamiento}{{A.6}{59}{Tema 6}{section*.92}{}}
\newlabel{sol:eje:6-desbordamiento@cref}{{[subappendix][6][2147483647,1]A.6}{[1][59][]59}}
\newlabel{sol:eje:6-redondeo}{{A.6}{60}{Tema 6}{section*.93}{}}
\newlabel{sol:eje:6-redondeo@cref}{{[subappendix][6][2147483647,1]A.6}{[1][60][]60}}
\makeatother

\begin{document}
\frontmatter
\begin{titlepage}
\centering
\vspace*{2.5cm}
{\large\scshape Anatomía de un LLM}\par
\vspace{0.4cm}
{\normalsize Fundamentos matemáticos de los modelos de lenguaje}\par
\vspace{2.2cm}
\rule{\textwidth}{0.8pt}\par\vspace{0.5cm}
{\small\scshape Fase 0 \textperiodcentered\ Preliminares matemáticos}\par\vspace{0.3cm}
{\Huge\bfseries Módulo 0.3\par}\vspace{0.3cm}
{\LARGE Números, álgebra y sumatorios\par}
\vspace{0.5cm}\rule{\textwidth}{0.4pt}\par
\vspace{1.6cm}
{\large Resumen del módulo}\par
\vfill
{\small Edición de octubre de 2026}\par
\end{titlepage}

\chapter*{Sobre este resumen}
\addcontentsline{toc}{chapter}{Sobre este resumen}
\markboth{Sobre este resumen}{}

Este documento reúne lo indispensable de los apuntes del Módulo 0.3: el texto expositivo, todas las definiciones y todos los resultados (teoremas, proposiciones, lemas y corolarios), las tablas y las observaciones de contenido. Se han suprimido las demostraciones, los ejemplos, las lecturas, las figuras, los ejercicios y sus soluciones. La numeración de temas, secciones, enunciados y ecuaciones es la misma que en los apuntes, de modo que cualquier referencia puede buscarse allí; las que remiten a elementos suprimidos (un ejemplo, una figura, una ecuación de una demostración) conservan su número en los apuntes.

\section*{Notación}
\begin{center}
\begin{tabular}{@{}ll@{}}
\toprule
Símbolo & Significado \\
\midrule
$\N=\{1,2,3,\dots\}$, $\N_0=\N\cup\{0\}$ & números naturales, con y sin el cero \\
$\Z$, $\Q$, $\R$ & enteros, racionales y reales \\
$\abs{x}$ & valor absoluto de $x$ \\
$\sum_{k=m}^{n}a_k$, $\prod_{k=m}^{n}a_k$ & sumatorio y productorio \\
$n!$, $\binom{n}{k}$ & factorial y número combinatorio \\
$\ln$, $\log_2$, $\log_{10}$ & logaritmos natural, binario y decimal \\
$\LSE(z)=\ln\sum_i e^{z_i}$ & función log-suma-exp \\
$(d_{m}\dots d_1d_0)_b$, \texttt{0x}\dots & número escrito en base $b$; prefijo hexadecimal \\
$\lfloor x\rfloor$, $\lceil x\rceil$ & parte entera por defecto y por exceso \\
$\fl(x)$, $\varepsilon$ & redondeo a coma flotante y épsilon de la máquina \\
\bottomrule
\end{tabular}
\end{center}


\tableofcontents

\mainmatter
% =====================================================================
\chapter{Los números reales}
\label{tema:1}
% =====================================================================

Este tema fija las reglas de juego de todo el módulo: qué se puede hacer con los números reales y por qué. En lugar de dar por sabidas las propiedades de la aritmética, se parte de una lista corta de axiomas, la de un \emph{cuerpo ordenado}, y se deducen de ella las reglas habituales. La presentación sigue los capítulos~1 y~2 de Spivak~\cite{spivak}; las secciones de potencias y radicales siguen los §1.1--1.3 de OpenStax~\cite{openstax}.

\section{Conjuntos numéricos}

Se trabaja con la cadena de inclusiones $\N\subset\Z\subset\Q\subset\R$, donde $\N=\{1,2,3,\dots\}$ son los naturales, $\Z$ los enteros, $\Q=\{p/q: p,q\in\Z,\ q\neq 0\}$ los racionales y $\R$ los reales. Un real que no es racional se llama \emph{irracional}.

\begin{proposicion}\label{prop:decimal-periodico}
Un número real es racional si y solo si su expresión decimal es finita o periódica.
\end{proposicion}

\addtocounter{definicion}{1}%

\section{Axiomas de cuerpo}

\begin{definicion}[Cuerpo]\label{def:cuerpo}
Un \emph{cuerpo} es un conjunto $K$ con dos operaciones, suma $(a,b)\mapsto a+b$ y producto $(a,b)\mapsto ab$, que cumplen, para todos $a,b,c\in K$:
\begin{enumerate}[label=(C\arabic*)]
  \item asociatividad: $(a+b)+c=a+(b+c)$ y $(ab)c=a(bc)$;
  \item conmutatividad: $a+b=b+a$ y $ab=ba$;
  \item elementos neutros: existen $0,1\in K$, con $0\neq 1$, tales que $a+0=a$ y $a\cdot 1=a$;
  \item opuesto: para cada $a$ existe $-a$ con $a+(-a)=0$;
  \item inverso: para cada $a\neq 0$ existe $a^{-1}$ con $aa^{-1}=1$;
  \item distributividad: $a(b+c)=ab+ac$.
\end{enumerate}
Se definen la resta $a-b=a+(-b)$ y, si $b\neq 0$, la división $a/b=ab^{-1}$.
\end{definicion}

$\Q$ y $\R$ son cuerpos; $\Z$ no lo es, porque $2$ no tiene inverso en $\Z$. Todas las reglas de la aritmética escolar se deducen de (C1)--(C6).

\begin{proposicion}[Consecuencias de los axiomas de cuerpo]\label{prop:cuerpo}
Para todos $a,b,x$ de un cuerpo:
\begin{enumerate}
  \item si $a+x=a$, entonces $x=0$ (unicidad del neutro); si $a+x=0$, entonces $x=-a$ (unicidad del opuesto);
  \item $a\cdot 0=0$;
  \item $(-a)b=-(ab)$ y $(-a)(-b)=ab$;
  \item si $ab=0$, entonces $a=0$ o $b=0$.
\end{enumerate}
\end{proposicion}

El apartado~(iv) es el que justifica resolver $(x-2)(x+3)=0$ igualando a cero cada factor.

\section{Axiomas de orden}

\begin{definicion}[Cuerpo ordenado]\label{def:orden}
Un cuerpo $K$ es \emph{ordenado} si existe un subconjunto $P\subset K$, el de los \emph{positivos}, tal que:
\begin{enumerate}[label=(O\arabic*)]
  \item tricotomía: para cada $a\in K$ se cumple exactamente una de $a=0$, $a\in P$, $-a\in P$;
  \item si $a,b\in P$, entonces $a+b\in P$;
  \item si $a,b\in P$, entonces $ab\in P$.
\end{enumerate}
Se escribe $a<b$ (o $b>a$) si $b-a\in P$, y $a\le b$ si $a<b$ o $a=b$.
\end{definicion}

\begin{proposicion}[Reglas de las desigualdades]\label{prop:orden}
En un cuerpo ordenado, para todos $a,b,c$:
\begin{enumerate}
  \item si $a<b$ y $b<c$, entonces $a<c$;
  \item si $a<b$, entonces $a+c<b+c$;
  \item si $a<b$ y $c>0$, entonces $ac<bc$;
  \item si $a\neq 0$, entonces $a^2>0$; en particular, $1>0$.
\end{enumerate}
\end{proposicion}

\begin{observacion}[La completitud]
$\Q$ y $\R$ son ambos cuerpos ordenados; lo que los distingue es el \emph{axioma del supremo}: todo subconjunto no vacío de $\R$ acotado superiormente tiene una cota superior mínima~\cite[cap.~8]{spivak}. Gracias a él existen en $\R$ números como $\sqrt 2$, que no están en $\Q$ (\cref{teo:raiz2}). En este módulo el axioma solo se usa implícitamente, al admitir que existen las raíces $n$-ésimas y los logaritmos de los reales positivos.
\end{observacion}

\section{Números irracionales}

\begin{lema}\label{lem:par}
Sea $n\in\Z$. Si $n^2$ es par, entonces $n$ es par.
\end{lema}

\begin{teorema}\label{teo:raiz2}
$\sqrt2$ es irracional.
\end{teorema}

\section{Potencias, radicales e intervalos}

Para $a\in\R$ y $n\in\N$ se define $a^n=a\cdots a$ ($n$ factores), $a^0=1$ si $a\neq 0$, y $a^{-n}=1/a^n$. Para $a>0$ se define $a^{1/n}=\sqrt[n]{a}$ como el único real positivo cuya potencia $n$-ésima es $a$, y $a^{m/n}=(a^{1/n})^m$.

\begin{proposicion}[Leyes de los exponentes]\label{prop:exponentes}
Para $a,b>0$ y $r,s\in\Q$: $a^{r}a^{s}=a^{r+s}$, $(a^{r})^{s}=a^{rs}$, $(ab)^{r}=a^{r}b^{r}$ y $a^{-r}=1/a^{r}$. Para exponentes enteros, las tres primeras valen también con $a,b\neq 0$.
\end{proposicion}

Para exponentes naturales, las leyes se prueban por inducción sobre el exponente; para exponentes racionales se reducen a ese caso usando la unicidad de la raíz $n$-ésima~\cite[§1.2--1.3]{openstax}.

\begin{definicion}[Intervalos]
Para $a<b$: $[a,b]=\{x: a\le x\le b\}$, $(a,b)=\{x: a<x<b\}$, y análogamente $[a,b)$, $(a,b]$, $(a,\infty)$, $(-\infty,b]$, etc. Los símbolos $\pm\infty$ no son números reales y nunca llevan corchete cerrado.
\end{definicion}

\addtocounter{definicion}{1}%

\addtocounter{definicion}{1}%


% =====================================================================
\chapter{Valor absoluto y desigualdades}
\label{tema:2}
% =====================================================================

Gran parte del trabajo matemático en aprendizaje automático consiste en \emph{acotar}: garantizar que un error es pequeño, que una pérdida no es negativa o que un método converge. Este tema reúne las desigualdades elementales que se usan para ello. El valor absoluto y la desigualdad triangular siguen el capítulo~1 de Spivak~\cite{spivak}; las inecuaciones, el §2.7 de OpenStax~\cite{openstax}; las demostraciones por inducción, el capítulo~10 de Hammack~\cite{hammack}.

\section{Valor absoluto}

\begin{definicion}[Valor absoluto]\label{def:abs}
Para $x\in\R$ se define $\abs{x}=x$ si $x\ge 0$ y $\abs{x}=-x$ si $x<0$. La \emph{distancia} entre $x$ e $y$ es $\abs{x-y}$.
\end{definicion}

\begin{proposicion}\label{prop:abs}
Para todos $x,y\in\R$ y $r\ge 0$:
\begin{enumerate}
  \item $\abs{x}\ge 0$, con igualdad si y solo si $x=0$; $\abs{-x}=\abs{x}$;
  \item $\abs{x}=\max(x,-x)$ y $-\abs{x}\le x\le\abs{x}$;
  \item $\abs{x}=\sqrt{x^2}$ y $\abs{xy}=\abs{x}\abs{y}$;
  \item $\abs{x}\le r$ si y solo si $-r\le x\le r$; $\abs{x}\ge r$ si y solo si $x\le -r$ o $x\ge r$.
\end{enumerate}
\end{proposicion}

\section{La desigualdad triangular}

\begin{teorema}[Desigualdad triangular]\label{teo:triangular}
Para todos $x,y\in\R$, $\abs{x+y}\le\abs{x}+\abs{y}$, con igualdad si y solo si $xy\ge0$.
\end{teorema}

\begin{corolario}[Desigualdad triangular inversa]\label{cor:triangular-inversa}
Para todos $x,y\in\R$, $\bigl|\abs{x}-\abs{y}\bigr|\le\abs{x-y}$.
\end{corolario}

La desigualdad triangular se generaliza a $n$ sumandos, $\abs{\sum_i x_i}\le\sum_i\abs{x_i}$ (\cref{eje:2-triangular-n}), y a vectores, donde $\abs{\cdot}$ se sustituye por una norma (módulo~2).

\section{Inecuaciones con valor absoluto}

Por la \cref{prop:abs}(iv), la condición $\abs{x-a}<r$ describe el intervalo $(a-r,a+r)$ de centro $a$ y radio $r$. Las inecuaciones con varios valores absolutos se resuelven separando casos según el signo de cada expresión.

\addtocounter{definicion}{1}%

\section{Desigualdad entre las medias aritmética y geométrica}

\begin{proposicion}[Caso de dos números]\label{prop:amgm2}
Si $a,b\ge0$, entonces $\sqrt{ab}\le\dfrac{a+b}{2}$, con igualdad si y solo si $a=b$.
\end{proposicion}

\begin{lema}\label{lem:producto1}
Si $x_1,\dots,x_n>0$ y $x_1x_2\cdots x_n=1$, entonces $x_1+\dots+x_n\ge n$, con igualdad si y solo si todos valen~1.
\end{lema}

\begin{teorema}[Desigualdad de las medias]\label{teo:amgm}
Si $a_1,\dots,a_n>0$, entonces
\[
\sqrt[n]{a_1a_2\cdots a_n}\le\frac{a_1+a_2+\dots+a_n}{n},
\]
con igualdad si y solo si $a_1=a_2=\dots=a_n$.
\end{teorema}

\addtocounter{definicion}{1}%

\section{Desigualdades de Bernoulli y de Cauchy--Schwarz}

\begin{teorema}[Desigualdad de Bernoulli]\label{teo:bernoulli}
Si $x\ge-1$ y $n\in\N$, entonces $(1+x)^n\ge1+nx$.
\end{teorema}

\begin{teorema}[Desigualdad de Cauchy--Schwarz]\label{teo:cs}
Para reales $a_1,\dots,a_n$ y $b_1,\dots,b_n$,
\[
\Bigl(\sum_{i=1}^na_ib_i\Bigr)^2\le\Bigl(\sum_{i=1}^na_i^2\Bigr)\Bigl(\sum_{i=1}^nb_i^2\Bigr).
\]
\end{teorema}

\addtocounter{definicion}{1}%

\addtocounter{definicion}{1}%


% =====================================================================
\chapter{Polinomios, ecuaciones y sistemas}
\label{tema:3}
% =====================================================================

Los polinomios son las funciones más sencillas después de las lineales y aparecen por todas partes: como aproximaciones de Taylor, como funciones de pérdida cuadráticas y como polinomio característico de una matriz. Este tema repasa su aritmética (división, raíces, factorización), la resolución de ecuaciones e inecuaciones y los sistemas lineales pequeños. Sigue los capítulos~2, 5 y~7 de OpenStax \emph{College Algebra}~\cite{openstax}, con las demostraciones que allí se omiten.

\section{Polinomios y división euclídea}

\begin{definicion}[Polinomio]
Un \emph{polinomio} con coeficientes reales es una expresión $P(x)=a_nx^n+\dots+a_1x+a_0$ con $a_i\in\R$. Si $a_n\neq0$, $n$ es el \emph{grado} de $P$, $\deg P$, y $a_n$ su \emph{coeficiente principal}. Un número $r$ es una \emph{raíz} de $P$ si $P(r)=0$.
\end{definicion}

Si $P$ y $Q$ son no nulos, $\deg(PQ)=\deg P+\deg Q$, porque el coeficiente principal del producto es el producto de los coeficientes principales, que no es cero (\cref{prop:cuerpo}(iv)).

\begin{teorema}[División euclídea]\label{teo:division}
Sean $P$ y $D$ polinomios con $D\neq0$. Existen polinomios únicos $Q$ (cociente) y $R$ (resto) tales que
\[
P=DQ+R\qquad\text{y}\qquad R=0\ \text{o}\ \deg R<\deg D.
\]
\end{teorema}

\begin{teorema}[Teoremas del resto y del factor]\label{teo:resto}
El resto de dividir $P$ entre $x-a$ es el número $P(a)$. En consecuencia, $a$ es raíz de $P$ si y solo si $x-a$ divide a $P$.
\end{teorema}

La \emph{regla de Ruffini} es una forma abreviada de dividir entre $x-a$: se bajan los coeficientes uno a uno, multiplicando el último resultado por $a$ y sumándolo al coeficiente siguiente.

\addtocounter{definicion}{1}%

\begin{corolario}\label{cor:raices}
Un polinomio no nulo de grado $n$ tiene a lo sumo $n$ raíces reales distintas.
\end{corolario}

\begin{teorema}[Raíces racionales]\label{teo:raices-racionales}
Sea $P(x)=a_nx^n+\dots+a_0$ con coeficientes enteros, $a_n\neq0$ y $a_0\neq0$. Si $p/q$, con $p,q$ enteros primos entre sí, es raíz de $P$, entonces $p$ divide a $a_0$ y $q$ divide a $a_n$.
\end{teorema}

\section{La ecuación de segundo grado y las fórmulas de Vieta}

\begin{teorema}[Ecuación de segundo grado]\label{teo:cuadratica}
Sean $a\neq0$ y $\Delta=b^2-4ac$ el \emph{discriminante}. La ecuación $ax^2+bx+c=0$ tiene dos soluciones reales distintas si $\Delta>0$, una (doble) si $\Delta=0$ y ninguna si $\Delta<0$. Cuando $\Delta\ge0$,
\[
x=\frac{-b\pm\sqrt{\Delta}}{2a}.
\]
\end{teorema}

\begin{proposicion}[Fórmulas de Vieta]\label{prop:vieta}
Si $x_1,x_2$ son las raíces (reales o repetidas) de $ax^2+bx+c$, entonces
\[
ax^2+bx+c=a(x-x_1)(x-x_2),\qquad x_1+x_2=-\frac ba,\qquad x_1x_2=\frac ca.
\]
Más en general, si $P(x)=a_n\prod_{i=1}^n(x-x_i)$, entonces $\sum_ix_i=-a_{n-1}/a_n$ y $\prod_ix_i=(-1)^na_0/a_n$.
\end{proposicion}

\addtocounter{definicion}{1}%

\section{Otras ecuaciones e inecuaciones}

Las ecuaciones \emph{bicuadradas} $ax^4+bx^2+c=0$ se reducen a segundo grado con $t=x^2$. En las ecuaciones \emph{con radicales} se aísla la raíz y se eleva al cuadrado; como elevar al cuadrado no es una operación reversible ($a^2=b^2$ no implica $a=b$), pueden aparecer soluciones \emph{extrañas} y es obligatorio comprobar cada candidata en la ecuación original. En las ecuaciones \emph{racionales} se multiplica por el denominador común y se descartan los valores que lo anulan.

Para resolver una inecuación $P(x)/Q(x)\ge0$ se factorizan numerador y denominador, se marcan en la recta los puntos donde se anula algún factor y se estudia el signo de cada factor en cada intervalo (\emph{tabla de signos}); los ceros del denominador nunca forman parte de la solución.

\addtocounter{definicion}{1}%

\section{Sistemas lineales pequeños}

Un sistema de $m$ ecuaciones lineales con $n$ incógnitas es \emph{compatible determinado} si tiene una única solución, \emph{compatible indeterminado} si tiene infinitas e \emph{incompatible} si no tiene ninguna. Los métodos de sustitución, igualación y reducción de bachillerato son casos particulares de la eliminación gaussiana del módulo~3.

\begin{teorema}[Regla de Cramer para $2\times2$]\label{teo:cramer}
El sistema $ax+by=e$, $cx+dy=f$ tiene solución única si y solo si $\delta=ad-bc\neq0$, y entonces
\[
x=\frac{ed-bf}{\delta},\qquad y=\frac{af-ec}{\delta}.
\]
\end{teorema}

\addtocounter{definicion}{1}%

\addtocounter{definicion}{1}%

\addtocounter{definicion}{1}%


% =====================================================================
\chapter{Sumatorios, productorios y binomio de Newton}
\label{tema:4}
% =====================================================================

La probabilidad de una secuencia es un producto de probabilidades condicionadas, su logaritmo es una suma, y la pérdida de entrenamiento es una media sobre millones de tokens. Manejar con soltura la notación $\sum$ y $\prod$ es, por tanto, imprescindible. Este tema sigue el capítulo~2 «Sums» de Graham, Knuth y Patashnik~\cite{gkp} y el capítulo~14 «Sums and Asymptotics» de Lehman, Leighton y Meyer~\cite{mcs}; el binomio de Newton sigue el §9.6 de OpenStax~\cite{openstax}.

\section{Notación y principio de inducción}

\begin{definicion}[Sumatorio y productorio]
Dados enteros $m\le n$ y números $a_m,\dots,a_n$,
\[
\sum_{k=m}^{n}a_k=a_m+a_{m+1}+\dots+a_n,\qquad \prod_{k=m}^{n}a_k=a_m\,a_{m+1}\cdots a_n.
\]
La letra $k$ es un \emph{índice mudo}: puede sustituirse por cualquier otra sin cambiar el valor. Si $n<m$, la suma es \emph{vacía} y vale~0, y el producto vacío vale~1. Más en general, $\sum_{k\in S}a_k$ denota la suma sobre un conjunto finito $S$ de índices. El \emph{factorial} es $n!=\prod_{k=1}^{n}k$, con $0!=1$.
\end{definicion}

Formalmente, la suma se define por recurrencia: $\sum_{k=m}^{m-1}a_k=0$ y $\sum_{k=m}^{n}a_k=\sum_{k=m}^{n-1}a_k+a_n$. Por eso las fórmulas para sumas se demuestran de forma natural por inducción.

\begin{teorema}[Principio de inducción]\label{teo:induccion}
Sea $P(n)$ una afirmación sobre los enteros $n\ge n_0$. Si $P(n_0)$ es cierta y, para cada $n\ge n_0$, $P(n)$ implica $P(n+1)$, entonces $P(n)$ es cierta para todo $n\ge n_0$.
\end{teorema}

Este principio es equivalente al principio del buen orden (todo subconjunto no vacío de $\N$ tiene mínimo); ambos se estudian en~\cite[cap.~2 y~5]{mcs} y~\cite[cap.~10]{hammack}.

\section{Propiedades de los sumatorios}

\begin{proposicion}\label{prop:sumas}
Para números $a_k,b_k$ y una constante $c$:
\begin{enumerate}
  \item linealidad: $\sum_{k=m}^n(ca_k+b_k)=c\sum_{k=m}^na_k+\sum_{k=m}^nb_k$;
  \item partición: si $m\le p<n$, $\sum_{k=m}^na_k=\sum_{k=m}^pa_k+\sum_{k=p+1}^na_k$;
  \item suma de una constante: $\sum_{k=m}^nc=(n-m+1)\,c$;
  \item cambio de índice: $\sum_{k=m}^na_k=\sum_{j=m+r}^{n+r}a_{j-r}$ para todo entero $r$, y $\sum_{k=0}^na_k=\sum_{k=0}^na_{n-k}$.
\end{enumerate}
\end{proposicion}

No hay reglas análogas para el producto: en general $\sum a_kb_k\neq\bigl(\sum a_k\bigr)\bigl(\sum b_k\bigr)$.

\section{Sumas clásicas}

\begin{teorema}[Suma aritmética]\label{teo:aritmetica}
Para todo $n\in\N$, $\displaystyle\sum_{k=1}^nk=\frac{n(n+1)}2$. En general, si $a_k=a_1+(k-1)d$, entonces $\sum_{k=1}^na_k=\dfrac{n(a_1+a_n)}{2}$.
\end{teorema}

\begin{teorema}[Suma geométrica]\label{teo:geometrica}
Para $r\neq1$ y $n\in\N_0$, $\displaystyle\sum_{k=0}^nr^k=\frac{1-r^{n+1}}{1-r}$. Si $r=1$, la suma vale $n+1$.
\end{teorema}

\begin{proposicion}[Sumas telescópicas]\label{prop:telescopica}
Para cualesquiera números $a_m,\dots,a_{n+1}$,
\[
\sum_{k=m}^n(a_{k+1}-a_k)=a_{n+1}-a_m.
\]
\end{proposicion}

\addtocounter{definicion}{1}%

\begin{teorema}[Sumas de potencias]\label{teo:potencias}
Para todo $n\in\N$,
\[
\sum_{k=1}^nk^2=\frac{n(n+1)(2n+1)}6,\qquad\sum_{k=1}^nk^3=\Bigl(\frac{n(n+1)}2\Bigr)^2.
\]
\end{teorema}

\addtocounter{definicion}{1}%

\section{Sumas dobles}

Una suma doble $\sum_{i}\sum_{j}a_{ij}$ suma los elementos de una tabla. Si los índices recorren un rectángulo, el orden de suma puede intercambiarse:
\[
\sum_{i=1}^m\sum_{j=1}^na_{ij}=\sum_{j=1}^n\sum_{i=1}^ma_{ij}.
\]
Si recorren un triángulo, hay que reescribir los límites~\cite[§2.4]{gkp}:
\[
\sum_{i=1}^n\sum_{j=i}^na_{ij}=\sum_{1\le i\le j\le n}a_{ij}=\sum_{j=1}^n\sum_{i=1}^ja_{ij}.
\]

\begin{proposicion}\label{prop:producto-sumas}
$\displaystyle\Bigl(\sum_{i=1}^ma_i\Bigr)\Bigl(\sum_{j=1}^nb_j\Bigr)=\sum_{i=1}^m\sum_{j=1}^na_ib_j$. En particular, $\displaystyle\Bigl(\sum_{i=1}^na_i\Bigr)^2=\sum_{i=1}^na_i^2+2\sum_{1\le i<j\le n}a_ia_j$.
\end{proposicion}

\addtocounter{definicion}{1}%

\section{Productorios y binomio de Newton}

Los productorios cumplen $\prod_k(a_kb_k)=\prod_ka_k\prod_kb_k$, $\prod_{k=1}^nc=c^n$ y, si todos los $a_k$ son positivos, $\ln\prod_ka_k=\sum_k\ln a_k$ (\cref{tema:5}).

\begin{definicion}[Números combinatorios]
Para enteros $0\le k\le n$, $\displaystyle\binom nk=\frac{n!}{k!\,(n-k)!}$, y $\binom nk=0$ si $k<0$ o $k>n$.
\end{definicion}

\begin{lema}[Identidad de Pascal]\label{lem:pascal}
Para $1\le k\le n$, $\displaystyle\binom nk+\binom n{k-1}=\binom{n+1}k$.
\end{lema}

\begin{teorema}[Binomio de Newton]\label{teo:binomio}
Para todos $a,b\in\R$ y $n\in\N_0$, $\displaystyle(a+b)^n=\sum_{k=0}^n\binom nka^{n-k}b^k$.
\end{teorema}

\begin{corolario}\label{cor:binomio}
$\displaystyle\sum_{k=0}^n\binom nk=2^n$ y, si $n\ge1$, $\displaystyle\sum_{k=0}^n(-1)^k\binom nk=0$.
\end{corolario}

\addtocounter{definicion}{1}%

\addtocounter{definicion}{1}%


% =====================================================================
\chapter{Exponenciales y logaritmos}
\label{tema:5}
% =====================================================================

El logaritmo convierte productos en sumas y potencias en productos. Esa propiedad lo hace omnipresente en el curso: la log-verosimilitud, la entropía, la perplejidad y la escala de las leyes de potencia son logaritmos. Este tema repasa las propiedades de exponenciales y logaritmos siguiendo el capítulo~6 de OpenStax~\cite{openstax}, introduce la función log-suma-exp siguiendo a Blanchard, Higham y Higham~\cite{blanchard} y fija las unidades de información (bits y nats) como en el capítulo~2 de Cover y Thomas~\cite{ct}. La construcción rigurosa de $e^x$ y $\ln x$ requiere el cálculo integral~\cite[cap.~18]{spivak}; aquí se admiten sus propiedades básicas.

\section{Potencias de exponente real y la función exponencial}

Para $a>0$, la potencia $a^r$ está definida para $r$ racional (\cref{tema:1}) y se extiende a todo $x\in\R$ de forma que la función $x\mapsto a^x$ es continua. Se admiten las siguientes propiedades, que extienden la \cref{prop:exponentes}.

\begin{proposicion}[Función exponencial]\label{prop:exp}
Sea $a>0$. Para todos $x,y\in\R$: $a^x>0$; $a^{x+y}=a^xa^y$; $(a^x)^y=a^{xy}$; $a^{-x}=1/a^x$. Si $a>1$, la función $x\mapsto a^x$ es estrictamente creciente; si $0<a<1$, es estrictamente decreciente. Si $a\neq1$, toma cada valor positivo exactamente una vez.
\end{proposicion}

El número $e=\lim_{n\to\infty}(1+\frac1n)^n\approx2{,}71828$ es la base natural; más en general, $e^x=\lim_{n\to\infty}(1+\frac xn)^n$ para todo $x$~\cite[cap.~18]{spivak}.

\begin{proposicion}\label{prop:ex-cota}
Para todo $x\in\R$, $e^x\ge1+x$.
\end{proposicion}

\section{Logaritmos}

\begin{definicion}[Logaritmo]\label{def:log}
Sean $a>0$, $a\neq1$ y $x>0$. El \emph{logaritmo en base $a$} de $x$ es el único $y\in\R$ tal que $a^y=x$; se escribe $y=\log_ax$. Se usan $\ln=\log_e$ (logaritmo natural), $\log_2$ (binario) y $\log_{10}$ (decimal).
\end{definicion}

Por definición, $a^{\log_ax}=x$ para $x>0$ y $\log_a(a^y)=y$ para todo $y$. El logaritmo solo está definido para argumentos positivos.

\begin{teorema}[Propiedades de los logaritmos]\label{teo:log}
Sean $a,b>0$ distintos de~1, $x,y>0$ y $r\in\R$. Entonces:
\begin{enumerate}
  \item $\log_a(xy)=\log_ax+\log_ay$ y $\log_a(x/y)=\log_ax-\log_ay$;
  \item $\log_a(x^r)=r\log_ax$;
  \item cambio de base: $\log_ax=\dfrac{\log_bx}{\log_ba}$; en particular, $\log_ab\cdot\log_ba=1$;
  \item si $a>1$, $x<y$ si y solo si $\log_ax<\log_ay$.
\end{enumerate}
\end{teorema}

\begin{corolario}\label{cor:lnx}
Para todo $x>0$, $\ln x\le x-1$, con igualdad si y solo si $x=1$. En consecuencia, $\ln x\ge1-\frac1x$.
\end{corolario}

\section{Ecuaciones exponenciales y logarítmicas}

Para resolver una ecuación exponencial se toman logaritmos o se cambia la incógnita ($t=a^x$); en las logarítmicas se usan las propiedades del \cref{teo:log} para llegar a $\log_aA=\log_aB$, que equivale a $A=B$, y se descartan las soluciones que hacen negativo o nulo algún argumento.

\addtocounter{definicion}{1}%

\addtocounter{definicion}{1}%

\section{La función log-suma-exp}

\begin{definicion}[Log-suma-exp]\label{def:lse}
Para $z=(z_1,\dots,z_n)\in\R^n$, $\displaystyle\LSE(z)=\ln\sum_{i=1}^ne^{z_i}$. La función \emph{softmax} es $\softmax(z)_i=e^{z_i}/\sum_je^{z_j}=e^{z_i-\LSE(z)}$.
\end{definicion}

\begin{proposicion}\label{prop:lse}
Para todo $z\in\R^n$ y $c\in\R$, con $M=\max_iz_i$:
\begin{enumerate}
  \item invariancia por traslación: $\LSE(z_1+c,\dots,z_n+c)=\LSE(z)+c$; en consecuencia, la softmax no cambia al sumar $c$ a todas las componentes;
  \item cotas: $M\le\LSE(z)\le M+\ln n$.
\end{enumerate}
\end{proposicion}

La \cref{prop:lse}(i) con $c=-M$ da el \emph{cálculo estable}: $\LSE(z)=M+\ln\sum_ie^{z_i-M}$. Todas las exponenciales son ahora $\le1$, de modo que no hay desbordamiento, y al menos una vale exactamente~1, de modo que el logaritmo no recibe un cero.

\addtocounter{definicion}{1}%

\section{Escalas logarítmicas, bits y nats}

Una cantidad que varía en muchos órdenes de magnitud se representa mejor por su logaritmo. Si $y=Cx^{k}$, entonces $\log y=\log C+k\log x$: en ejes doblemente logarítmicos una \emph{ley de potencia} es una recta de pendiente $k$, que es como se presentan las leyes de escala del módulo~16. Si $y=Ca^{x}$, entonces $\log y$ es lineal en $x$ (escala semilogarítmica).

\begin{definicion}[Unidades de información]
La \emph{información} de un suceso de probabilidad $p\in(0,1]$ es $-\log_2p$ \emph{bits} o $-\ln p$ \emph{nats}~\cite[cap.~2]{ct}. Por el cambio de base, $1\ \text{nat}=\frac1{\ln2}\ \text{bits}\approx1{,}4427$ bits.
\end{definicion}

\addtocounter{definicion}{1}%

\addtocounter{definicion}{1}%

\addtocounter{definicion}{1}%


% =====================================================================
\chapter{Sistemas de numeración y coma flotante}
\label{tema:6}
% =====================================================================

Un ordenador solo almacena bits. Este tema estudia cómo se representan con ellos los enteros sin signo, los enteros con signo y los números reales, y qué se pierde por el camino. La representación posicional se fundamenta en la división euclídea de enteros; la parte de máquina sigue el capítulo~2 «Representing and Manipulating Information» de Bryant y O'Hallaron~\cite{csapp} y el artículo de Goldberg~\cite{goldberg}; los formatos de coma flotante son los de la norma IEEE~754~\cite{ieee754}.

\section{Representación posicional}

\begin{teorema}[Representación en base $b$]\label{teo:base}
Sea $b\ge2$ un entero. Todo $n\in\N$ se escribe de manera única como
\[
n=\sum_{k=0}^{m}d_kb^k=(d_md_{m-1}\dots d_1d_0)_b,\qquad d_k\in\{0,1,\dots,b-1\},\ d_m\neq0.
\]
Los $d_k$ son las \emph{cifras} de $n$ en base $b$.
\end{teorema}

La demostración es un algoritmo: las cifras de $n$ en base $b$ son los restos de dividir sucesivamente entre $b$, leídos de abajo arriba. En sentido contrario, el valor de $(d_m\dots d_0)_b$ se obtiene con el \emph{esquema de Horner}, $(\cdots((d_mb+d_{m-1})b+d_{m-2})\cdots)b+d_0$, que es la regla de Ruffini del \cref{tema:3} evaluando en $x=b$.

\addtocounter{definicion}{1}%

\begin{proposicion}[Número de cifras]\label{prop:cifras}
Un entero $n\ge1$ tiene exactamente $\lfloor\log_bn\rfloor+1$ cifras en base $b$. Con $m$ cifras en base $b$ se representan exactamente los $b^m$ enteros $0,1,\dots,b^m-1$.
\end{proposicion}

Las cifras de la parte fraccionaria de un número $x\in[0,1)$ se obtienen multiplicando sucesivamente por $b$ y tomando la parte entera: $x=\sum_{k\ge1}d_kb^{-k}=(0{,}d_1d_2\dots)_b$.

\begin{proposicion}\label{prop:fraccion-binaria}
Sea $p/q$ una fracción irreducible con $0<p<q$. Su desarrollo en base~2 es finito si y solo si $q$ es una potencia de~2.
\end{proposicion}

Por ejemplo, $0{,}1=1/10$ tiene denominador $10=2\cdot5$, de modo que su desarrollo binario es infinito: $0{,}1=(0{,}0\overline{0011})_2$. Ningún ordenador binario almacena $0{,}1$ exactamente.

\section{Binario, hexadecimal y bytes}

Un \emph{bit} es una cifra binaria y un \emph{byte} son 8 bits, que representan $2^8=256$ valores. Como $16=2^4$, cada cifra \emph{hexadecimal} ($0$--$9$, $\texttt{A}$--$\texttt{F}$) corresponde a un grupo de 4 bits, y un byte se escribe con dos cifras hexadecimales, de $\texttt{0x00}$ a $\texttt{0xFF}$. Para convertir entre binario y hexadecimal basta agrupar los bits de cuatro en cuatro desde la derecha: $(1011\,0110)_2=\texttt{0xB6}=182$.

Para cantidades de memoria se usan dos sistemas de prefijos: los decimales del SI ($1\ \text{kB}=10^3$ bytes, $1\ \text{GB}=10^9$ bytes) y los binarios de la IEC ($1\ \text{KiB}=2^{10}$ bytes, $1\ \text{GiB}=2^{30}$ bytes). La diferencia es de un $7{,}4\,\%$ en el caso de los gigas.

\section{Enteros con signo: complemento a dos}

\begin{definicion}[Complemento a dos]\label{def:c2}
Una palabra de $n$ bits $w=(a_{n-1}\dots a_0)$ representa, \emph{sin signo}, el entero $U(w)=\sum_{k=0}^{n-1}a_k2^k$ y, en \emph{complemento a dos}, el entero
\[
T(w)=-a_{n-1}2^{n-1}+\sum_{k=0}^{n-2}a_k2^k=U(w)-a_{n-1}2^n.
\]
\end{definicion}

\begin{teorema}\label{teo:c2}
Para $n\ge2$:
\begin{enumerate}
  \item $T$ es una biyección entre las palabras de $n$ bits y los enteros del intervalo $[-2^{n-1},\,2^{n-1}-1]$; $T(w)<0$ si y solo si $a_{n-1}=1$;
  \item $T(w)$ y $U(w)$ difieren en un múltiplo de $2^n$;
  \item si $\bar w$ es la palabra con todos los bits invertidos, $T(\bar w)+1=-T(w)$; en particular, el opuesto de $T(w)$ se obtiene invirtiendo los bits y sumando~1, y el resultado es representable en $n$ bits salvo si $T(w)=-2^{n-1}$.
\end{enumerate}
\end{teorema}

\addtocounter{definicion}{1}%

Por el \cref{teo:c2}(ii), la suma de enteros con signo se realiza con el mismo circuito que la suma sin signo, descartando el acarreo final: el resultado es correcto salvo que salga del rango, situación que se llama \emph{desbordamiento} (\cref{eje:6-desbordamiento}).

\section{Coma flotante}

\begin{definicion}[Número en coma flotante normalizado]\label{def:flotante}
Un formato binario de coma flotante con $e$ bits de exponente, sesgo $\beta=2^{e-1}-1$ y $p$ bits de fracción representa la palabra $(s,E,f)$, con $1\le E\le2^e-2$, como
\[
x=(-1)^s\,(1{,}f_1f_2\dots f_p)_2\cdot2^{E-\beta}.
\]
Los valores extremos de $E$ se reservan para el cero, los números subnormales, los infinitos y \texttt{NaN}~\cite{ieee754}.
\end{definicion}

\begin{table}[h]
\centering\small
\caption{Formatos binarios de coma flotante habituales. El épsilon es $\varepsilon=2^{-p}$, la distancia de~1 al siguiente número representable.}\label{tab:formatos}
\begin{tabular}{@{}lcccccc@{}}
\toprule
Formato & Bits & Signo & Exponente $e$ & Fracción $p$ & $\varepsilon=2^{-p}$ & Máximo finito \\
\midrule
fp64 (doble) & 64 & 1 & 11 & 52 & $2{,}2\cdot10^{-16}$ & $1{,}8\cdot10^{308}$ \\
fp32 (simple) & 32 & 1 & 8 & 23 & $1{,}2\cdot10^{-7}$ & $3{,}4\cdot10^{38}$ \\
fp16 (media) & 16 & 1 & 5 & 10 & $9{,}8\cdot10^{-4}$ & $65\,504$ \\
bf16 (bfloat16) & 16 & 1 & 8 & 7 & $7{,}8\cdot10^{-3}$ & $3{,}4\cdot10^{38}$ \\
\bottomrule
\end{tabular}
\end{table}

El formato bfloat16 conserva los 8 bits de exponente de fp32 y recorta la fracción a 7 bits: tiene el mismo rango que fp32 con mucha menos precisión, y sus 16 bits son exactamente los 16 bits superiores de un fp32.

\begin{proposicion}[Error de redondeo]\label{prop:redondeo}
Sea $\fl(x)$ el número del formato más cercano a $x$. Si $x$ está en el rango de los números normalizados, entonces
\[
\fl(x)=x(1+\delta)\qquad\text{con}\qquad\abs{\delta}\le u=2^{-(p+1)}=\frac\varepsilon2.
\]
\end{proposicion}

La demostración es el \cref{eje:6-redondeo}. El número $u$ se llama \emph{unidad de redondeo}: el error \emph{relativo} es pequeño, pero el absoluto crece con $\abs{x}$, porque la distancia entre números consecutivos en $[2^k,2^{k+1})$ es $2^{k-p}$~\cite{goldberg}.

\addtocounter{definicion}{1}%

\addtocounter{definicion}{1}%



\backmatter
% =====================================================================
\begin{thebibliography}{99}
\addcontentsline{toc}{chapter}{Bibliografía}
\markboth{Bibliografía}{}
\small

\bibitem{openstax}
J.~Abramson.
\newblock \emph{College Algebra}, 2.ª ed.
\newblock OpenStax, Rice University, 2021.
\newblock \url{https://openstax.org/details/books/college-algebra-2e}

\bibitem{blanchard}
P.~Blanchard, D.~J. Higham y N.~J. Higham.
\newblock Accurate computation of the log-sum-exp and softmax functions.
\newblock Preprint, 2019.
\newblock \url{https://arxiv.org/abs/1909.03469}

\bibitem{csapp}
R.~E. Bryant y D.~R. O'Hallaron.
\newblock \emph{Computer Systems: A Programmer's Perspective}, 3.ª ed.
\newblock Pearson, 2016.
\newblock Cap.~2, «Representing and Manipulating Information». Texto de los cursos de sistemas de Carnegie Mellon.

\bibitem{ct}
T.~M. Cover y J.~A. Thomas.
\newblock \emph{Elements of Information Theory}, 2.ª ed.
\newblock Wiley, 2006.

\bibitem{goldberg}
D.~Goldberg.
\newblock What every computer scientist should know about floating-point arithmetic.
\newblock \emph{ACM Computing Surveys}, 23(1):5--48, 1991.

\bibitem{gkp}
R.~L. Graham, D.~E. Knuth y O.~Patashnik.
\newblock \emph{Concrete Mathematics: A Foundation for Computer Science}, 2.ª ed.
\newblock Addison-Wesley, 1994.
\newblock Cap.~2, «Sums».

\bibitem{hammack}
R.~Hammack.
\newblock \emph{Book of Proof}, 3.ª ed.
\newblock Richmond (Virginia), 2018.
\newblock \url{https://richardhammack.github.io/BookOfProof/}

\bibitem{ieee754}
IEEE.
\newblock \emph{IEEE Standard for Floating-Point Arithmetic} (IEEE Std 754-2019).
\newblock IEEE, 2019.

\bibitem{mcs}
E.~Lehman, F.~T. Leighton y A.~R. Meyer.
\newblock \emph{Mathematics for Computer Science}, revisión del 6 de junio de 2018.
\newblock Massachusetts Institute of Technology. Texto del curso 6.042J.
\newblock \url{https://courses.csail.mit.edu/6.042/spring18/mcs.pdf}

\bibitem{mit6042}
A.~R. Meyer y A.~Chlipala.
\newblock 6.042J: Mathematics for Computer Science.
\newblock Massachusetts Institute of Technology, MIT OpenCourseWare, primavera de 2015.
\newblock \url{https://ocw.mit.edu/courses/6-042j-mathematics-for-computer-science-spring-2015/}

\bibitem{spivak}
M.~Spivak.
\newblock \emph{Calculus}, 4.ª ed.
\newblock Publish or Perish, 2008.
\newblock Cap.~1, «Basic Properties of Numbers», y cap.~2, «Numbers of Various Sorts».

\end{thebibliography}


\end{document}
